%==============================================================================
%  CHAPTER 5 — SPRINT 2: TRAINING PIPELINES & EXPERIMENT TRACKING
%==============================================================================
\chapter{Sprint 2 --- Training Pipelines and Experiment Tracking}
\label{chap:sprint2}

\section*{Introduction}
The second sprint delivers the heart of the platform: uploading code and data,
running it as an isolated training job on the cluster, and automatically capturing
the results. This chapter follows the same structure as the previous one.

% ---------------------------------------------------------------------------
\section{Sprint Backlog}
\begin{table}[H]\centering
\caption{Sprint~2 backlog}\label{tab:backlog-s2}
\begin{tabularx}{\textwidth}{p{0.6cm} X p{2cm}}
\toprule
\textbf{\#} & \textbf{Task} & \textbf{Pts}\\
\midrule
3 & Upload code/dataset to object storage & 3\\
3 & Auto-detect the entry script in an upload & 2\\
4 & Build the training pipeline (download $\to$ run $\to$ autolog) & 5\\
4 & Convert notebooks to scripts and inject the autolog wrapper & 3\\
5 & Enable automatic capture of metrics, parameters, artifacts & 3\\
6 & Stream live pipeline logs to the UI & 3\\
7 & Build the experiments browsing interface & 3\\
\bottomrule
\end{tabularx}
\end{table}

% ---------------------------------------------------------------------------
\section{Analysis and Conception}

\subsection{Use Cases}
This sprint covers \emph{Upload code and data}, \emph{Trigger a training run},
\emph{Watch live logs}, and \emph{Browse experiment runs}.

\subsection{Sequence Diagram: Triggering a Training Run}
% >>> USER: render diagrams/seq_training.puml -> images/seq_training.pdf
\begin{figure}[H]\centering
  % \includegraphics[width=\textwidth]{seq_training.pdf}
  \fbox{\parbox[c][5.5cm][c]{\textwidth}{\centering\itshape
    Placeholder --- insert \texttt{images/seq\_training.pdf}\\
    (render \texttt{diagrams/seq\_training.puml})}}
  \caption{Sequence diagram of triggering a custom-code training run}
  \label{fig:seq-training}
\end{figure}

\subsection{Data Model Increment}
The sprint adds the \texttt{PipelineRun} entity (pipeline run identifier, status,
parameters, timestamps) linked to a project.

% ---------------------------------------------------------------------------
\section{Realization}

\subsection{Upload and Storage}
Uploaded files are stored in the object store under a key derived from the project
and a short unique prefix. The backend supports plain scripts, notebooks, and
archives, and can auto-detect the entry script when one is not specified, following
a fixed priority order.

\subsection{The Training Pipeline}
A training run is expressed as a small pipeline that executes inside an isolated
pod: it downloads the code (and dataset) from the object store, installs any
missing dependencies, and runs the entry script. Notebooks are first converted to
a Python script. The script is not run directly, it is wrapped by a generated
runner that enables automatic experiment tracking before handing control to the
user's code.

\begin{lstlisting}[language=Python,caption={The generated runner enables autologging before running the user's script.},label={lst:runner}]
import os, runpy, mlflow
mlflow.set_tracking_uri(os.environ["MLFLOW_TRACKING_URI"])
mlflow.set_experiment(experiment_id=os.environ["MLFLOW_EXPERIMENT_ID"])
mlflow.autolog(log_models=True)          # one line instruments sklearn/XGBoost/...
runpy.run_path("user_script.py", run_name="__main__")
\end{lstlisting}

\subsection{Robust Notebook Handling}
Because real-world notebooks are rarely written for headless execution, the runner
includes several robustness measures: it strips interactive notebook magics during
conversion, neutralizes blocking calls, and provides compatibility shims for
runtime-only modules so that code exported from cloud notebook environments runs
unchanged.

\subsection{Live Logs}
While a run executes, the backend reads the pod's logs from the cluster and streams
them to the interface, so the user watches training progress in real time. The
\texttt{autolog} call records metrics, parameters, and artifacts to the tracking
server with no changes to the user's code.

\subsection{User Interface}
% >>> USER: screenshots: upload view, live-log terminal, experiments list.
The interface offers a drag-and-drop upload area, a live log terminal, and an
experiments view listing each run with its captured metrics.

% ---------------------------------------------------------------------------
\section{Sprint Review and Retrospective}
At the end of Sprint~2 a user can upload arbitrary code, run it on the cluster, and
see metrics and artifacts captured automatically, the platform now \emph{trains
and tracks}.
\begin{figure}[H]\centering
\begin{tikzpicture}
\begin{axis}[
    width=0.8\textwidth, height=6cm,
    xlabel={Day}, ylabel={Remaining points},
    xmin=0, xmax=12, ymin=0, ymax=24,
    legend pos=north east, grid=both,
]
\addplot[dashed,gray] coordinates {(0,22)(11,0)};
\addplot[thick,primary,mark=*] coordinates {
  (0,22)(1,22)(2,20)(3,18)(4,16)(5,13)(6,11)(7,9)(8,7)(9,4)(10,2)(11,0)};
\legend{Ideal, Actual}
\end{axis}
\end{tikzpicture}
\caption{Sprint~2 burndown chart}\label{fig:burndown-s2}
\end{figure}

\noindent\textbf{Retrospective.} % >>> USER
The training and tracking pipeline was the most technically demanding increment;
handling messy real-world notebooks required extra robustness work. \todo{your note}.
